\section{Theoretische Grundlagen}
\label{sec:theorie}

Die theoretischen Grundlagen des Versuchs, inklusive der fundamentalen Gleichungen.

\subsection{Grundlegende Gesetze}
Die zentrale Beziehung des Versuchs ist z. B. das Newtonsche Gesetz
\begin{equation}
  F = m \cdot a \label{eq:newton}
\end{equation}
mit Masse $m$ in \si{\kilogram} und Beschleunigung $a$ in
\si{\meter\per\square\second}.

\subsection{Wichtige Formeln}
Weitere Formeln können hier folgen, z. B. eine Winkelabhängigkeit:
\begin{equation}
  U(\theta) = U_0 \cdot \cos^2(\theta)
  \label{eq:cos2}
\end{equation}

Für die Fehlerrechnung nutzen wir die Gauß'sche Fehlerfortpflanzung:
\begin{equation}
  \gausserror
  \label{eq:gauss}
\end{equation}

Abbildungen zur Theorie, z. B. ein Diagramm oder eine Skizze:
\begin{figure}[H]
  \centering
  % \includegraphics[width=0.6\textwidth]{theorie_skizze.pdf}
  \caption{Platzhalter: Skizze zur Theorie (\cref{eq:cos2}).}
  \label{fig:theorie_skizze}
\end{figure}
